\documentclass[11pt,oneside]{book}
\usepackage[
    paperwidth=6in,
    paperheight=9in,
    top=0.75in,
    bottom=0.75in,
    left=0.6in,
    right=0.6in,
    twoside=false
]{geometry}
\input{../pri-end/T0_preamble_De.tex}
% T0_preamble_patches.tex — LOKALER PLATZHALTER
% Im Repo durch die offizielle Version ersetzen.

\usepackage{graphicx}

\makeatletter
\renewcommand{\thechapter}{\arabic{chapter}}
\makeatother
\raggedbottom
\allowdisplaybreaks
\hfuzz=5pt
\vfuzz=5pt
\emergencystretch=2em

\setcounter{tocdepth}{0}

\begin{document}
\frontmatter
\pagestyle{fancy}

% ── Titelseite ──────────────────────────────────────────────
\begin{titlepage}
  \centering\vspace*{1.5cm}
  {\Huge\bfseries\color{blue!70!black}
    Dunkle Materie — eine Illusion}\\[0.6cm]
  {\Large\bfseries\color{blue!40!black}
    Warum Galaxien ohne unsichtbare Masse funktionieren}\\[1cm]
  {\LARGE\itshape Fundamentale Fraktale Geometrische\\[0.2cm]
    Feldtheorie (FFGFT) — Galaktische Physik}\\[0.5cm]
  {\large\itshape Ein populärwissenschaftlicher Überblick}\\[2cm]
  {\large Johann Pascher}\\[0.3cm]
  {\normalsize ORCID: 0009-0000-6518-4064}\\[0.3cm]
  {\normalsize\url{https://github.com/jpascher/T0-Time-Mass-Duality}}\\[1.5cm]
  {\normalsize September 2026}
  \vfill{\small Version 1.0 --- Populärwissenschaftliche Ausgabe}
\end{titlepage}

\newpage\thispagestyle{empty}%
{\small\noindent\textbf{Stichworte:} Dunkle Materie $\cdot$
Rotationskurven $\cdot$ Galaxien $\cdot$ MOND $\cdot$
Gravitationslinsen $\cdot$ Bullet Cluster $\cdot$ FFGFT $\cdot$
Trägheitsregime\\[1em]
\noindent Johann Pascher: \textit{Dunkle Materie — eine Illusion.
Warum Galaxien ohne unsichtbare Masse funktionieren. FFGFT.}\\
Erste Ausgabe, September 2026.\\[0.6em]
\noindent Lizenz: Creative Commons BY 4.0}

\tableofcontents
\mainmatter

% ════════════════════════════════════════════════════════════
\chapter{Das Rätsel der flachen Kurven}
% ════════════════════════════════════════════════════════════

Es war eine einfache Messung mit einer unbequemen Antwort.

In den 1970er Jahren beobachtete Vera Rubin etwas, das die
Physik seither nicht losgelassen hat: die Sterne am Rand einer
Galaxie bewegen sich zu schnell. Viel zu schnell. Nach dem
Gesetz der Schwerkraft — demselben, das Planeten um die Sonne
kreisen lässt — müssten weit außen liegende Sterne langsamer
sein als die weiter innen. So wie Neptun langsamer um die Sonne
kreist als die Erde.

Aber das tun sie nicht.

Die Rotationskurven von Galaxien — also die Kurve, die zeigt,
wie schnell Sterne in verschiedenen Abständen vom Zentrum
kreisen — werden außen flach. Die Geschwindigkeit bleibt
konstant, so weit man messen kann. Das ist, als würden alle
Planeten in unserem Sonnensystem mit derselben Geschwindigkeit
die Sonne umkreisen, egal wie weit sie entfernt sind.

Die Standarderklärung lautet: Es muss mehr Masse geben, als
wir sehen. Eine unsichtbare, leuchtlose Masse — Dunkle Materie
— muss die Galaxien umhüllen und die Gravitation verstärken.
Ohne sie würden die Sterne am Rand wegfliegen.

Diese Erklärung klingt vernünftig. Sie hat aber einen Preis:
Dunkle Materie wurde nie direkt nachgewiesen. Weder in
Teilchenbeschleunigern noch in unterirdischen Detektoren noch
durch irgendein anderes Experiment. Wir postulieren 26~\%
des gesamten Universums aus einem Stoff, den wir nicht kennen,
nicht sehen und nicht greifen können.

Dieses Buch stellt eine andere Frage: Was wenn die Galaxien
sich nicht falsch verhalten — sondern unsere Physik an ihren
Rändern aufhört zu gelten?

Die FFGFT — Fundamentale Fraktale Geometrische Feldtheorie —
gibt eine präzise Antwort darauf. Sie braucht keine neuen
Teilchen. Sie braucht keine unsichtbare Masse. Sie braucht
nur eine einzige Zahl: $\xi \approx \tfrac{4}{3} \times 10^{-4}$.

% ════════════════════════════════════════════════════════════
\chapter{Was Newton wirklich sagt — und wo er aufhört}
% ════════════════════════════════════════════════════════════

Newtons Gravitationsgesetz ist eine der verlässlichsten
Formeln der Physik. Auf der Erde, im Sonnensystem, bei der
Berechnung von Satellitenbahnen — es funktioniert mit
außerordentlicher Präzision.

Aber es hat eine versteckte Annahme: Es setzt voraus, dass
Masse und Trägheit unter allen Umständen dieselbe Beziehung
haben. Trägheit ist das, was einen Körper daran hindert, seine
Geschwindigkeit zu ändern — je größer die Trägheit, desto mehr
Kraft braucht man, um ihn zu beschleunigen.

Newton schrieb: $F = m \cdot a$. Kraft gleich Masse mal
Beschleunigung. Das klingt selbstverständlich. Aber was wenn
dieser Zusammenhang nur für große Beschleunigungen gilt?

In unserem Alltag sind Beschleunigungen groß. Ein fallendes
Objekt auf der Erde erfährt $g \approx 10~\text{m/s}^2$.
Ein Satellit im Erdorbit vielleicht noch ein Tausendstel davon.
Aber ein Stern am äußersten Rand einer Galaxie erfährt eine
Beschleunigung von nur
$a \sim 10^{-10}~\text{m/s}^2$ — das ist ein
Zehnmilliardtel der Erdbeschleunigung.

Haben wir jemals geprüft, ob Newton in diesem Regime gilt?

Die ehrliche Antwort ist: Nein. Im Labor, im Sonnensystem, auf
der Erde — dort haben wir sein Gesetz tausendfach bestätigt.
Im galaktischen Tiefenregime haben wir es nie direkt getestet.

Die FFGFT sagt: Dort gilt eine andere Physik. Nicht weil eine
neue Kraft auftaucht, sondern weil die Trägheit selbst
nichtlinear wird.

% ════════════════════════════════════════════════════════════
\chapter{Die Grundrelation: Zeit und Masse als Kehrwerte}
% ════════════════════════════════════════════════════════════

Der Kern der FFGFT ist eine einzige Formel:

\[
  \tilde{T} \cdot m = 1 .
\]

Intrinsische Zeit und Masse sind Kehrwerte voneinander.
Wo Masse groß ist, vergeht Zeit langsam. Wo Masse klein ist,
vergeht Zeit schnell. Das ist nicht metaphorisch gemeint —
es ist eine geometrische Eigenschaft des Feldes.

Aus dieser Relation folgt: Ein Körper mit sehr kleiner Masse
lebt in einem Zeitregime, das sich fundamental von dem eines
schweren Körpers unterscheidet. Und da Trägheit mit Masse und
Zeit zusammenhängt, verändert sich auch die Trägheit.

Das klingt abstrakt. Aber es hat eine ganz konkrete Konsequenz
für Galaxien.

Am Rand einer Galaxie ist die Masse der sichtbaren Materie
gering, und die Beschleunigung sehr klein. Dort treten beide
Effekte zusammen auf: kleine Masse und kleine Beschleunigung.
Die FFGFT sagt, dass in diesem Regime die Newton'sche Formel
$F = ma$ nicht mehr exakt gilt. Die Trägheit wird kleiner als
Newton voraussagt — Sterne reagieren auf Schwerkraft
empfindlicher als erwartet.

Das ist der Effekt, den wir in den Rotationskurven sehen —
ohne Dunkle Materie.

% ════════════════════════════════════════════════════════════
\chapter{Die Übergangsskala $a_0$ — eine Grenze aus erster Ableitung}
% ════════════════════════════════════════════════════════════

Damit eine Theorie wissenschaftlich ist, muss sie präzise
vorhersagen, wo ein Effekt beginnt. Die FFGFT tut das.

Es gibt eine Beschleunigungsskala $a_0$, unterhalb derer das
Newton'sche Regime endet und das galaktische Regime beginnt.
Diese Skala ist keine freie Zahl, die man anpasst — sie folgt
direkt aus dem einzigen Parameter $\xi$ der Theorie:

\[
  a_0 = \frac{c^2 \xi^{10}}{4\,\bar{\lambda}_e}
    \approx 1{,}03 \times 10^{-10}~\text{m/s}^2 .
\]

Dabei ist $\bar{\lambda}_e = \hbar/(m_e c)$ die reduzierte
Compton-Wellenlänge des Elektrons — eine bekannte, messbare
Größe. Der Faktor $\xi^{10}$ trägt die gesamte
Skalenstruktur: zehn Potenzen der Grundkonstante verbinden
die Quantenskala des Elektrons mit der Skala galaktischer
Beschleunigung.

Das ist bemerkenswert. Denn dieser Wert $a_0$ war in der
Astrophysik seit Jahrzehnten empirisch bekannt — Mordehai
Milgrom hatte ihn 1983 aus Beobachtungen bestimmt, ohne ihn
herleiten zu können. Die FFGFT leitet ihn aus erster
Ableitung her, ohne freien Parameter.

Unterhalb von $a_0$ gilt das modifizierte Kraftgesetz:

\[
  a = \sqrt{a_N^2 + a_N\,a_0} ,
\]

wobei $a_N$ die Newton'sche Beschleunigung aus der sichtbaren
Masse ist. Für große Beschleunigungen ($a_N \gg a_0$) geht
dieser Ausdruck in Newton über: $a \to a_N$. Für kleine
Beschleunigungen ($a_N \ll a_0$) ergibt sich:

\[
  a \to \sqrt{a_N\,a_0} \qquad\Longrightarrow\qquad
  v^4 = G\,M_b\,a_0 .
\]

Das ist die \emph{baryonische Tully-Fisher-Relation} — eine
der verlässlichsten empirischen Regeln der Galaktikenphysik,
die seit vierzig Jahren beobachtet wird, ohne dass das
Standardmodell sie je aus Grundprinzipien hergeleitet hätte.
Die FFGFT leitet sie her.

% ════════════════════════════════════════════════════════════
\chapter{Der Test: DDO 154 und die Milchstraße}
% ════════════════════════════════════════════════════════════

Eine Theorie muss sich an Beobachtungen messen lassen. Drei
Galaxien unterschiedlichen Typs sind besonders aufschlussreich.

\textbf{DDO 154 — der sauberste Test.}
Diese Zwerggalaxie ist gasdominiert: fast ihre gesamte
sichtbare Masse besteht aus Wasserstoffgas, dessen Menge
radioastronomisch sehr genau bekannt ist. Es gibt kaum
Unsicherheit über die sichtbare Masse — das macht DDO 154
zum idealen Prüfstein.

Die Newton'sche Vorhersage für die äußere Rotationsgeschwindigkeit
beträgt rund 18,5~km/s. Die Beobachtung zeigt 47,0~km/s —
mehr als doppelt so viel.

Die FFGFT-Vorhersage aus dem modifizierten Kraftgesetz: 47,1~km/s.

Das Verhältnis Vorhersage zu Beobachtung: 1,001 — eine
Übereinstimmung auf einem Promille.

\textbf{NGC 3198 — eine Spiralgalaxie.}
Hier ist die baryonische Masse weniger genau bekannt (Unsicherheit
$\sim 30\,\%$). Die Newton'sche Erwartung: 50,4~km/s,
beobachtet: 150~km/s. FFGFT sagt 121~km/s voraus — eine
Abweichung von 18\,\%, die vollständig im Bereich der
Massenunicherheit liegt.

\textbf{Die Milchstraße.}
Newton erwartet 122~km/s in der äußeren Scheibe, gemessen
werden rund 220~km/s, FFGFT sagt 180~km/s. Auch hier liegt
die Abweichung im Unsicherheitsbereich der Sternenmasse.

\medskip
Dunkle Materie erklärt diese Kurven ebenfalls — aber nur,
indem man für jede Galaxie ein eigenes Haloprofil anpasst.
Das sind freie Parameter, einer pro Galaxie. Die FFGFT hat
keinen einzigen freien Parameter: $a_0$ folgt aus $\xi$,
das selbst die gesamte Physik trägt.

% ════════════════════════════════════════════════════════════
\chapter{Der Bullet Cluster — das scheinbar stärkste Argument}
% ════════════════════════════════════════════════════════════

Wer gegen Dunkle Materie argumentiert, bekommt früher oder
später denselben Einwand: den Bullet Cluster.

2006 veröffentlichten Douglas Clowe und Kollegen eine
Beobachtung, die als entscheidender Beweis für Dunkle Materie
gilt. Zwei Galaxiencluster sind miteinander kollidiert. Das
heiße Gas — die weitaus größte sichtbare Massenkomponente —
wurde durch die Kollision abgebremst und blieb zurück. Die
Galaxien selbst flogen durch.

Aus Gravitationslinsen-Messungen bestimmte man, wo die
Masse tatsächlich sitzt: nicht beim Gas, sondern bei den
Galaxien. Das Gas wurde abgebremst, die Masse aber nicht.
Ergo: Es muss dort eine kollisionsfreie unsichtbare Masse
geben — Dunkle Materie.

Das ist ein überzeugendes Argument. Aber die FFGFT beantwortet
es ohne Dunkle Materie.

Das modifizierte Kraftgesetz wirkt auf alle Masse — sichtbare
und interagierende gleichermaßen. Entscheidend ist das
Massenverhältnis zwischen Galaxien und Gas in einem
Galaxiencluster. Die Galaxien tragen rund zehnmal mehr Masse
als das Gas. Bei der Kollision wird das Gas durch den
Strahlungsdruck und elektromagnetische Wechselwirkung
abgebremst — die Galaxien nicht.

In einer Finite-Elemente-Simulation der Kollisionsgeometrie
(mit den Parametern von Clowe et al.\ 2006) ergibt die FFGFT:
Der Gravitationslinsen-Peak liegt 10~kpc von den Galaxien
entfernt und 240~kpc vom Gaszentrum. Genau das wird beobachtet.

Der Bullet Cluster ist kein Beweis für Dunkle Materie. Er ist
ein Beweis dafür, dass Masse und Gas bei einer Kollision
anders reagieren. Das war schon immer bekannt — und folgt
in der FFGFT direkt aus der Geometrie.

% ════════════════════════════════════════════════════════════
\chapter{Dunkle Energie — ein Konstrukt der falschen Lesart}
% ════════════════════════════════════════════════════════════

Neben Dunkler Materie kennt das Standardmodell eine zweite
Unbekannte: Dunkle Energie. Sie soll die beschleunigte
Expansion des Universums antreiben und macht angeblich 69\,\%
des gesamten Universums aus.

Zusammen belaufen sich die beiden unbekannten Komponenten auf
95\,\% des Universums. Nur 5\,\% sind normale Materie, die
wir kennen und verstehen.

Die FFGFT braucht weder Dunkle Materie noch Dunkle Energie —
aber aus unterschiedlichen Gründen.

Dunkle Materie ist, wie gezeigt, kein Stoff: Es ist ein Regime,
in dem die Trägheit nichtlinear wird. Die Effekte sind real —
die Rotationskurven sind real, die Gravitationslinsen sind real
— aber sie brauchen keine neue Zutat.

Dunkle Energie verschwindet aus einem anderen Grund: Es gibt
keine Expansion, die zu beschleunigen wäre. Das Universum ist
in der FFGFT statisch. Die Rotverschiebung weit entfernter
Galaxien entsteht nicht weil der Raum sich dehnt, sondern
weil Licht auf seinem Weg durch das $\xi$-Feld Energie verliert:

\[
  1 + z = e^{\xi\,x} .
\]

Für kleine Abstände ist das ununterscheidbar vom Doppler-Effekt.
Für große Abstände macht es dieselben Vorhersagen wie das
Standardmodell — weil die Daten entartet sind: Sie können
prinzipiell nicht zwischen Expansion und geometrischer
Wegverlängerung unterscheiden.

Die kosmologische Konstante $\Lambda$ — das mathematische
Kürzel für Dunkle Energie — hat in einem statischen Universum
keinen Referenten. Sie ist eine Buchungsgröße der
Expansions-Lesart, kein physikalisches Objekt.

% ════════════════════════════════════════════════════════════
\chapter{Was die Daten wirklich entscheiden können}
% ════════════════════════════════════════════════════════════

Die Frage ist berechtigt: Wenn die FFGFT dieselben Beobachtungen
erklären kann wie das Standardmodell — warum sollte man dann
wechseln?

Die Antwort liegt in der Struktur der Theorien, nicht in
einzelnen Datenpunkten.

Das Standardmodell der Kosmologie ($\Lambda$CDM) beschreibt
die Kosmologie mit sechs freien Parametern, die alle empirisch
bestimmt werden: Hubble-Konstante, Baryonendichte, Dunkle-Materie-Dichte,
Spektralindex, Fluktuationsamplitude, optische Tiefe. Dazu
kommen die 19 Parameter des Teilchenphysik-Standardmodells.
Insgesamt über 25 Zahlen, von denen keine aus Grundprinzipien
hergeleitet ist.

Die FFGFT hat einen einzigen Parameter: $\xi = \tfrac{4}{3}
\times 10^{-4}$. Aus ihm folgen:
\begin{itemize}
  \item die Gravitationskonstante $G$,
  \item die Hubble-Konstante $H_0 = 66{,}82~\text{km/s/Mpc}$,
  \item die Beschleunigungsskala $a_0 = 1{,}03 \times 10^{-10}~\text{m/s}^2$,
  \item die CMB-Temperatur $T_\text{CMB} = 2{,}725~\text{K}$,
  \item die baryonische Tully-Fisher-Relation,
  \item die Leptonmassen-Hierarchie.
\end{itemize}

Derselbe $\xi$ trägt fünf verschiedene Sektoren der Physik —
Teilchenphysik, Kosmologie, galaktische Dynamik, Gravitationslinsen,
Quantenmechanik. Das ist keine Kalibrierung: Es ist
Querverankerung.

Die kosmologische Konstante $\Lambda$ dagegen arbeitet nirgendwo
außer im eigenen Fit. Das ist der Unterschied zwischen einer
Theorie und einer Buchungsgröße.

\medskip

Wo kann man heute entscheiden? Die klassischen Säulen —
Supernovae, baryonische akustische Oszillationen, CMB-Spektrum
— sind entartet: Beide Modelle liefern dieselben
mathematischen Kurven. Aber es gibt Beobachtungen, die diskriminieren:

\begin{itemize}
  \item \textbf{Wellenlängenabhängigkeit der Rotverschiebung:}
    FFGFT sagt eine streng achromatische Rotverschiebung voraus.
    Jede nachgewiesene Wellenlängenabhängigkeit würde das Modell
    falsifizieren.
  \item \textbf{Direkte Dunkle-Materie-Detektion:}
    Jeder direkte Nachweis eines Dunkle-Materie-Teilchens —
    in Beschleunigern, Detektoren, kosmischen Quellen — würde
    zeigen, dass tatsächlich ein neuer Stoff existiert.
  \item \textbf{Euclid und DESI:}
    Die neuen Galaxienvermessungen werden das Wachstum
    kosmischer Strukturen mit bisher unerreichter Präzision
    messen. Die FFGFT macht konkrete Vorhersagen für das
    Leistungsspektrum ohne Dunkle-Materie-Halo-Korrekturen.
\end{itemize}

% ════════════════════════════════════════════════════════════
\chapter{Eine ehrliche Bilanz}
% ════════════════════════════════════════════════════════════

Die FFGFT erklärt Galaxienrotationskurven ohne Dunkle Materie.
Sie leitet die Beschleunigungsskala $a_0$ aus erster Ableitung her.
Sie besteht den Bullet-Cluster-Test. Sie erklärt, warum Dunkle
Energie kein physikalisches Objekt ist.

Das ist keine vollständige Theorie der Gravitation. Es gibt
offene Fragen:

\textbf{Der K3-Vorwärtsbeweis} — die mathematische Herleitung
der Geometrie, aus der das modifizierte Kraftgesetz folgt,
ist strukturell begründet, aber noch nicht vollständig bewiesen.

\textbf{14\,\% Abweichung bei $a_0$:} Der theoretisch abgeleitete
Wert liegt knapp 14\,\% über der Mitte der empirischen Spanne.
Die Literaturwerte für $a_0$ streuen allerdings zwischen
$1{,}1$ und $1{,}3 \times 10^{-10}~\text{m/s}^2$ — der
FFGFT-Wert liegt innerhalb dieser Spanne.

\textbf{Clusterphysik:} Bei großen Galaxienclustern mit
bekannter Massenunicherheit sind die Vorhersagen weniger
präzise als bei gasdominerten Zwerggalaxien.

\medskip

Was aber klar ist: Die Annahme, dass 95\,\% des Universums
aus unbekannten Stoffen bestehen, die nie direkt nachgewiesen
wurden, ist keine kleinere Annahme als die, dass die Trägheit
bei sehr kleinen Beschleunigungen nichtlinear wird.

Der Unterschied: Die eine Annahme ist nicht falsifizierbar
— man kann immer mehr Dunkle Materie hinzufügen, wenn eine
Beobachtung nicht passt. Die andere hat eine präzise,
ausgezeichnete Skala $a_0$, die aus der Struktur der Theorie
folgt und sich in Beobachtungen zeigt.

Dunkle Materie ist nicht falsch. Sie könnte existieren. Aber
sie ist nicht notwendig. Und eine Physik, die ohne sie
auskommt, sollte ernsthaft geprüft werden.

% ════════════════════════════════════════════════════════════
\chapter{Was weitergeht}
% ════════════════════════════════════════════════════════════

Die nächste Generation von Beobachtungen wird entscheiden.

Das Euclid-Teleskop der ESA vermesst Milliarden von Galaxien
bis in Rotverschiebungen von $z \sim 2$. Das James-Webb-Teleskop
zeigt uns Galaxien in einem frühen Universum, die nach dem
Standardmodell eigentlich noch nicht existieren dürften —
aber im statischen FFGFT-Universum problemlos Platz haben.
DESI misst baryonische akustische Oszillationen mit
Rekordpräzision.

Für die FFGFT sind konkrete Vorhersagen offen:

Die achromatische Rotverschiebung — keine Wellenlängenabhängigkeit
über den gesamten beobachtbaren Bereich — ist ein Kerntest.
Aktuelle Präzisionsbeobachtungen setzen bereits enge Grenzen;
die FFGFT ist bisher kompatibel.

Die Massenverteilung in Galaxienclustern ohne Dunkle-Materie-Halo:
Die Clusterprofile der FFGFT unterscheiden sich subtil von
NFW-Profilen der Standard-Dunkle-Materie-Modelle. Die
Auflösung zukünftiger Linsen-Surveys könnte das unterscheiden.

Die baryonische Tully-Fisher-Relation gilt nach FFGFT exakt
— nicht als statistische Tendenz mit Streuung, sondern als
geometrische Folge. Abweichungen wären ein direktes Indiz
für Dunkle Materie als echten Stoff.

\medskip

Die Frage, ob Dunkle Materie existiert oder eine Illusion
der falschen Physik ist, ist eine der größten offenen Fragen
der Wissenschaft. Sie ist entscheidbar. Die Theorie, die sie
beantwortet, wird nicht durch Glauben gewonnen — sondern durch
Vorhersagen, die sich entweder bewähren oder scheitern.

Die FFGFT steht bereit, gemessen zu werden.

\backmatter

% ── Quellen ──────────────────────────────────────────────────
\chapter*{Quellen und Grundlagendokumente}
\addcontentsline{toc}{chapter}{Quellen}

Dieses Buch basiert auf dem FFGFT-Korpus im GitHub-Repository
\url{https://github.com/jpascher/T0-Time-Mass-Duality}.
Die folgenden Dokumente sind die primären Quellen:

\begin{itemize}
  \item \textbf{Dok.~308} — Der kosmische Sektor der FFGFT
    (Rotverschiebung, Dunkle Materie, Bullet Cluster, Lichtablenkung)
  \item \textbf{Dok.~267} — Kosmologische Entartung:
    $\Lambda$CDM und FFGFT als zwei Lesarten derselben Beobachtungen
  \item \textbf{Dok.~309} — Das Skalenanker-Problem
    ($H_0$-Herleitung, SI-Anker, Asymmetrie der Modelle)
  \item \textbf{Dok.~025} — T0-Kosmologie
    (CMB, $\xi$-Feld-Effekte, statisches Universum)
  \item \textbf{Dok.~004} — T0-Modell Übersicht
    (Grundrelationen, Feldgeometrien)
  \item \textbf{Dok.~012} — Gravitationskonstante aus FFGFT
\end{itemize}

Alle Berechnungen in Kapitel~4 und~5 wurden mit den
Verifikationsskripten
\texttt{ffgft\_308\_p44\_stufeA\_a0\_kette.py} und
\texttt{ffgft\_308\_p44\_stufeC\_bullet\_cluster.py}
überprüft.

\end{document}
